\scanpage{130}
(par ailleurs arbitraire). Nous supposerons $L_\bullet\in\operatorname{Ob}D^-(A)$, et dirons que $L_\bullet$ est ``pseudo-cohérent'' s'il est isomorphe dans $D^-(A)$ à un complexe qui est \uline{libre de type fini} en tout degré; \inserted{si $A$ est noeth., c'est équivalent au fait que les $H^i(L_\bullet)$ de type fini.} \inserted{On dit} qu'il est de ``Tor-dimension finie'' s'il est isomorphe dans $D^-(A)$ à un complexe \uline{plat} qui est \deleted{\unsure{de type}} à degrés bornés. On voit aisément que ces conditions sont stables par \unsure{mapping cylindre} et par $\overset{\mathbb L}{\otimes}$.

[\unsure{donc|Tous} les complexes satisfaisant l'une ou l'autre de ces conditions font \inserted{encore} une catégorie triangulée]. Les complexes pseudo-cohérents de Tor-dimension finie sont donc exactement ceux qui sont isomorphes dans $D^b(A)$ \deleted{\unsure{aux}} à des complexes qui sont \uline{projectifs de type fini} en $\forall$ degré, et qui sont \uline{de plus à degrés bornés}.

Si $\inserted{u_\bullet}\deleted{\unsure{f}}$ est un endomorphisme d'un tel $L_\bullet$, on peut lui associer \inserted{de façon unique} un
\mspage{129}
